Consider a realistic problem: a company provides calculation services from various servers to clients, who have requirements whether they be server performance or location. The company, to maximize revenue, may represent the servers and clients as two anticliques of a bipartite graph, the server eligibility for client as edges and solve the resulting Bipartite Matching problem to service as many clients as possible. If the company doesn't know what clients will request services then the problem is Online Bipartite Matching instead.
This problem is well-studied, with a known optimal algorithm, and has received significant attention for the simple reason that a bipartite graph can easily represent a consumer-producer relation. A large amount of effort has been given to this problem due to it's relation to advertisement, as it can simulate a simple advertising model with individual ads and impressions forming bipartite anticliques and the further possibility to generalize algorithms to the AdWords problem which fully simulates current advertising strategies
There is, however, an issue regarding using Online Bipartite Matching as a model for real problems, that being that it is often too restrictive and algorithms based on it aren't as efficient as they could be. In some problems the algorithm may be able to wait, or make small corrections which could improve the size of the found matching but cannot be reflected in the Online Bipartite Matching problem.
In this work we study the performance of various algorithms in Online Bipartite Matching and related models that either increase the capabilities of the algorithm, like allowing for a longer wait time before matching, or restricting the adversary, such as forcing the graph to be revealed in a random order